\documentclass[11pt]{article}

% --- Page geometry ---
\usepackage[margin=1in]{geometry}

% --- Mathematics ---
\usepackage{amsmath}
\usepackage{amssymb}
\usepackage{amsthm}

% --- Graphics and tables ---
\usepackage{graphicx}
\usepackage{booktabs}

% --- Cross-referencing ---
\usepackage{hyperref}
\usepackage[capitalise,noabbrev]{cleveref}

% --- Bibliography ---
\usepackage{natbib}
\bibliographystyle{plainnat}

% --- Miscellaneous ---
\usepackage{url}

% --- Theorem environments ---
\newtheorem{theorem}{Theorem}[section]
\newtheorem{proposition}[theorem]{Proposition}
\newtheorem{lemma}[theorem]{Lemma}
\theoremstyle{definition}
\newtheorem{definition}[theorem]{Definition}
\theoremstyle{remark}
\newtheorem{observation}[theorem]{Observation}
\newtheorem{remark}[theorem]{Remark}

% --- Notation shortcuts ---
\newcommand{\R}{\mathbb{R}}
\newcommand{\abs}[1]{\lvert #1 \rvert}
\newcommand{\norm}[1]{\lVert #1 \rVert}
\newcommand{\Pe}{\mathrm{Pe}}
\newcommand{\dt}{\Delta t}
\newcommand{\dx}{\Delta x}
\DeclareMathOperator{\coth}{coth}

% ============================================================
\title{Nonstandard Finite Difference Schemes for the Black--Scholes Equation:\\
Implementation and Numerical Validation in QuantLib}
\author{}
\date{}

\begin{document}
\maketitle

\begin{abstract}
We present the implementation and numerical validation of three spatial
discretization schemes for the one-dimensional Black--Scholes partial
differential equation in log-price coordinates: a standard central-difference
scheme, the exponential-fitting scheme of \citet{Duffy2004}, and an adapted
Milev--Tagliani scheme based on the CN-effective-diffusion construction
of \citet{MilevTagliani2010a,MilevTagliani2010b}. All three schemes are
realized within the finite-difference framework of a QuantLib fork, sharing a
common mesher, operator assembly, and time-stepping infrastructure. We derive
tridiagonal stencil coefficients in log-space, prove adapted positivity and
stability results, and validate the implementation through eight numerical
experiment groups encompassing truncated-call oscillation tests, discrete
double-barrier pricing with Monte Carlo benchmarks, grid convergence studies,
effective-diffusion and M-matrix diagnostics, performance benchmarks, and
P\'eclet-number regime visualization. The results confirm that
exponential fitting provides unconditional M-matrix compliance at moderate
overhead, while the Milev--Tagliani adaptation offers competitive accuracy
in the low-volatility regime under Crank--Nicolson time stepping, with a
safe fallback to exponential fitting when its prerequisites are violated.
\end{abstract}

% ============================================================
\section{Introduction}
\label{sec:intro}

% [task3 + task20: Introduction and scope/caveats]

% ============================================================
\section{Mathematical Framework}
\label{sec:math}

% ------------------------------------------------------------
\subsection{The Black--Scholes PDE}
\label{sec:bspde}

% [task4: BS PDE in S-space and log-space]

% ------------------------------------------------------------
\subsection{Spatial Discretization Schemes}
\label{sec:schemes}

% [task4: Three scheme derivations with tridiagonal coefficients]

\subsubsection{Standard Central Differences}
\label{sec:sc}

\subsubsection{Exponential Fitting}
\label{sec:ef}

\subsubsection{Milev--Tagliani CN-Effective Diffusion}
\label{sec:mt}

% ------------------------------------------------------------
\subsection{P\'eclet Number and Regime Analysis}
\label{sec:peclet}

% [task5: Péclet definition, xCothx evaluation]

% ------------------------------------------------------------
\subsection{M-Matrix Property and Positivity}
\label{sec:mmatrix}

% [task5: M-matrix definition, EF guarantee proof]

% ------------------------------------------------------------
\subsection{Adapted Positivity and Stability Results}
\label{sec:theorems}

% [task6: Adapted Theorems 3.1 and 3.2]

% ------------------------------------------------------------
\subsection{Time Discretization and CN-Equivalence}
\label{sec:time}

% [task7: CN, CN-equivalence in 1D, dampingSteps]

% ============================================================
\section{Implementation in QuantLib}
\label{sec:impl}

% [task9: Architecture, CN-equivalence gate, fallback, barriers]

% ============================================================
\section{Numerical Experiments}
\label{sec:experiments}

% [task10: experiment descriptions; task11: figures; task12-15: tables]

% ============================================================
\section{Results and Discussion}
\label{sec:results}

% [task16: key findings, recommendations, literature comparison]

% ============================================================
\section{Conclusion}
\label{sec:conclusion}

% [task17: contributions, recommendations, future work]

% ============================================================
\appendix

\section{Adapted Proofs of Positivity and Stability}
\label{app:proofs}

% [task18: Full proofs of adapted Theorems 3.1 and 3.2]

\section{Benchmark Timing Details}
\label{app:benchmark}

% [task19: Detailed timing tables with machine specs]

% ============================================================
\bibliography{refs}

\end{document}
